% !TeX TS-program = xelatex
\documentclass{resume-photo}

\usepackage{xeCJK}
\setCJKmainfont{Noto Serif CJK SC}
\usepackage{fontspec}
\usepackage{geometry}
\geometry{top=0.55cm, bottom=0.45cm, left=0.85cm, right=0.85cm, includefoot}
\usepackage{enumitem}
\setlist[itemize]{itemsep=0pt, topsep=1pt, parsep=0pt, partopsep=0pt}
\linespread{1.05}

\ResumeName{陈晓明}

\begin{document}

\ResumeContacts{
  187-XXXX-8821,%
  \ResumeUrl{mailto:chenxm@sjtu.edu.cn}{chenxm@sjtu.edu.cn},%
  \textnormal{上海交通大学 | 软件工程 · 本科 | 2021-XX}%
}

\ResumeTitle

% ====================== 简介 ======================
\noindent {\small 上海交通大学软件工程本科毕业，3 年前端与全栈开发经验。熟练使用 React / TypeScript 构建复杂交互应用，深度参与 Node.js 服务端开发与性能优化，有 ToC 大流量产品从 0 到 1 落地经验。注重工程质量与用户体验，具备独立负责完整功能模块的能力。}

% ====================== 教育 ======================
\section{教育背景}
\ResumeItem{上海交通大学}[\textnormal{软件学院 | 软件工程 | 本科 | 专业排名前 10\%}][2018.09—2022.06]

% ====================== 技能 ======================
\section{专业技能}
\begin{itemize}
  \item \textbf{前端}：React / Next.js / Vue 3 / TypeScript；熟悉 Webpack / Vite 构建优化；掌握 CSS-in-JS、Tailwind CSS。
  \item \textbf{后端}：Node.js / Express / NestJS；了解 Golang；熟悉 RESTful API 与 GraphQL 设计。
  \item \textbf{工程}：熟悉 CI/CD（GitHub Actions / Jenkins）、Docker、Nginx；有 AWS / 阿里云部署经验。
  \item \textbf{其他}：PostgreSQL / Redis / MongoDB；熟悉性能监控（Sentry / Lighthouse）；了解微前端架构。
\end{itemize}

% ====================== 工作 ======================
\section{工作经历}

\ResumeItem{字节跳动｜飞书前端团队}[\textnormal{高级前端工程师}][2023.07—至今]
\begin{itemize}
  \item 负责飞书文档编辑器核心功能迭代，基于 Slate.js 实现协同编辑冲突解决算法，将协同冲突率从 0.8\% 降至 0.1\%。
  \item 主导编辑器渲染性能优化：引入虚拟滚动与增量渲染，万行文档首屏时间从 3.2s 降至 0.9s，LCP 指标提升 70\%。
  \item 设计并落地组件库规范，抽象 40+ 通用组件，覆盖 3 个业务团队，开发效率提升 30\%。
  \item 参与飞书文档移动端适配，实现 PC/移动双端统一 Slate 模型，减少 40\% 冗余代码。
\end{itemize}

\ResumeItem{美团｜到店前端团队}[\textnormal{前端工程师}][2022.07—2023.06]
\begin{itemize}
  \item 从零搭建商家端营销活动配置平台（React + NestJS），支持 10 万+ 商家每日创建活动，系统可用性 99.9\%。
  \item 实现低代码页面搭建引擎，拖拽配置生成活动落地页，商家配置时间从平均 2 小时缩短至 15 分钟。
  \item 优化 SSR 渲染链路，活动页 TTFB 降低 60\%，Google PageSpeed 得分从 58 提升至 91。
\end{itemize}

% ====================== 项目 ======================
\section{项目经历}

\ResumeItem{开源项目：react-virtual-grid}[][2023.01—至今]
\begin{itemize}
  \item 基于 React 18 开发高性能虚拟化表格组件，支持百万行数据渲染，内存占用相比 ag-Grid 降低 55\%。
  \item GitHub 获得 1.2k+ stars，npm 周下载量 8000+，已被 5 个开源项目引用。
\end{itemize}

\ResumeItem{个人全栈项目：AI 写作助手}[][2022.10—2022.12]
\begin{itemize}
  \item 基于 Next.js 13 + OpenAI API 实现流式输出写作助手，集成 Supabase 实现用户系统与文档存储。
  \item 上线 3 个月获得 500+ 注册用户，日活 50+，Vercel 部署 P95 响应时间 < 200ms。
\end{itemize}

% ====================== 荣誉 ======================
\section{个人荣誉}
\begin{itemize}
  \item 字节跳动 2024 年度"最具影响力工程师"（全飞书前端 Top 5\%）
  \item ACM-ICPC 亚洲区域赛铜奖；上海交通大学 2022 年度优秀毕业生
  \item 掘金技术社区认证作者，技术文章累计阅读量 30 万+；GitHub 个人主页 500+ followers
\end{itemize}

\end{document}
